% ============================================================
% Paragraph: Synthetic Data Calibration Results
% For use in the Experiments section of the thesis/paper
% ============================================================

%--- Simulation Setup ---%
To evaluate the calibration performance of the proposed method in a controlled setting, a synthetic option chain is constructed from a single set of known Heston parameters: $\kappa = 2.0$, $\theta = 0.1$, $\sigma = 0.4$, $\rho = -0.6$, $v_0 = 0.05$, $r = 0.05$, and $S_0 = 1000$. These parameters satisfy the Feller condition ($2\kappa\theta = 0.40 > \sigma^2 = 0.16$), ensuring that the variance process remains strictly positive. A total of 99 valid call option prices are generated via QuantLib across a grid of five maturities $\tau \in \{0.3, 0.4, 0.5, 0.6, 0.7\}$ years and 20 equally spaced moneyness levels $K/S_0 \in [0.6, 1.5]$. Market implied volatilities and Vega values are subsequently computed from the synthetic prices using the Black--Scholes inversion. Both methods under comparison---the proposed DDN with IV-space calibration and the baseline method of Zhang~(2025)---are applied to this identical synthetic chain using multi-start Adam optimisation with 10 restarts, a learning rate of $5\times10^{-3}$, and a maximum of 500 gradient steps per restart. Full details of the experimental setup are provided in Table~\ref{tab:synthetic_setup}.

%--- Parameter Recovery ---%
Table~\ref{tab:param_recovery} reports the recovered Heston parameters and their relative errors with respect to the ground truth for both methods. Neither method achieves accurate point-wise parameter recovery for all five parameters, which is expected given the well-known non-identifiability of the Heston model: multiple parameter combinations can produce nearly identical option price surfaces. Notably, both methods substantially overestimate $\kappa$ (by approximately 91\% and 79\%, respectively) while recovering $\rho$ and $\theta$ with considerably lower errors (below 16\%). The proposed method recovers $\sigma$ with a relative error of 36.61\%, compared to 78.49\% for Zhang~(2025), demonstrating a material advantage in volatility-of-variance estimation. Calibration times are comparable for the two methods (11.76\,s vs.\ 11.30\,s), confirming that the added Feller regularisation and IV-space objective do not introduce significant computational overhead. Importantly, the Feller constraint is satisfied by the parameters returned by the proposed method ($2\kappa\theta = 0.644 > \sigma^2 = 0.299$), whereas Zhang~(2025) does not enforce this physical constraint.

%--- Re-pricing Accuracy ---%
Table~\ref{tab:repricing_accuracy} summarises the re-pricing accuracy on the synthetic chain. The proposed method achieves a mean IV relative error of 3.08\% and a maximum of 9.25\% across all 99 contracts, with an IV MAE of 0.00781. In price space, the QuantLib re-pricing with the calibrated parameters yields a mean relative price error of 12.78\%. These errors are primarily driven by deep out-of-the-money (OTM) contracts ($K/S_0 \gtrsim 1.3$), where near-zero option prices cause small absolute deviations to inflate relative metrics. For Zhang~(2025), the mean price re-pricing error is 184.80\% with a maximum of 7873.29\%, and only 63.6\% of contracts are priced within a 5\% relative error. The substantially lower price-space errors of the proposed method (12.78\% vs.\ 184.80\%) confirm that calibrating in IV space with Vega weighting and enforcing the Feller condition yields parameter estimates that are more physically consistent and produce better aggregate re-pricing performance, even though both methods exhibit the inherent parameter degeneracy characteristic of the Heston model.
